\documentclass[11pt,a4paper]{article}

\usepackage[utf8]{inputenc}
\usepackage[T1]{fontenc}
\usepackage{lmodern}
\usepackage[margin=2.4cm]{geometry}
\usepackage{array}
\usepackage{longtable}
\usepackage{hyperref}
\usepackage{xcolor}

\hypersetup{
  colorlinks=true,
  linkcolor=blue,
  urlcolor=blue,
  pdftitle={Using Claude Code With OpenCode Go as the Backend},
  pdfauthor={claude-o-go}
}

\setlength{\parindent}{0pt}
\setlength{\parskip}{0.7em}
\renewcommand{\arraystretch}{1.2}

\newcommand{\code}[1]{\texttt{#1}}

\title{\Huge Using Claude Code With OpenCode Go\\[0.4em]\Large A practical backend tutorial for \code{claude-o-go}}
\author{\code{claude-o-go}}
\date{\today}

\begin{document}
\maketitle
\tableofcontents
\newpage

\section{Goal}

This tutorial explains the setup in this repository and why it works.
The goal is to keep the Claude Code CLI experience while sending model
requests to OpenCode Go instead of Anthropic's Claude API.

Claude Code remains the harness: terminal UI, tool calling, file edits,
permissions, prompt handling, and the agent loop. OpenCode Go is the backend
that receives the actual model request.

\section{What You Have}

This directory contains:

\begin{itemize}
  \item \code{.env}: your OpenCode Go API key.
  \item \code{claude-o-go}: the launcher script you run instead of \code{claude}.
  \item \code{opencode-go-claude-proxy.mjs}: the local proxy between Claude Code and OpenCode Go.
  \item \code{opencode-go-claude-example.mjs}: a direct API example that calls OpenCode Go without Claude Code.
\end{itemize}

The normal command is:

\begin{verbatim}
cd /home/user/Desktop/opencode-go-w-claude
./claude-o-go
\end{verbatim}

For a one-shot prompt:

\begin{verbatim}
./claude-o-go -p "Reply with exactly: OK"
\end{verbatim}

\section{Run \code{claude-o-go} From Anywhere}

You do not need to change into this repository every time, and you do not need
to copy it into each new project. Install it once and use \code{claude-o-go}
from any project directory.

\subsection{One-Time Setup}

Make the launcher executable if needed:

\begin{verbatim}
chmod +x /home/user/Desktop/opencode-go-w-claude/claude-o-go
\end{verbatim}

Create a symlink in a directory already on your \code{PATH}:

\begin{verbatim}
mkdir -p ~/.local/bin
ln -s /home/user/Desktop/opencode-go-w-claude/claude-o-go \
  ~/.local/bin/claude-o-go
\end{verbatim}

Make sure \code{\textasciitilde{}/.local/bin} is on your \code{PATH}:

\begin{verbatim}
export PATH="$HOME/.local/bin:$PATH"
\end{verbatim}

Then reload your shell:

\begin{verbatim}
source ~/.bashrc   # or: source ~/.zshrc
\end{verbatim}

\subsection{Why This Works}

The launcher resolves its own symlink with \code{readlink -f}. It then changes
into the real install directory, where \code{.env} and
\code{opencode-go-claude-proxy.mjs} live. That makes the launcher work no
matter where you call it from.

\subsection{Optional Shell Profile Configuration}

If you would rather not keep the key in \code{.env}, export it once in your
shell profile:

\begin{verbatim}
export OPENCODE_GO_API_KEY="sk-..."
export OPENCODE_GO_MODEL="qwen3.7-plus"
\end{verbatim}

The launcher picks up \code{OPENCODE\_GO\_API\_KEY} from the environment and
skips reading \code{.env}. The model is no longer hardcoded in the launcher;
define \code{OPENCODE\_GO\_MODEL} in your shell profile so it applies
everywhere.

\subsection{Daily Usage}

From any project directory:

\begin{verbatim}
cd /path/to/any/project
claude-o-go
\end{verbatim}

For a one-shot:

\begin{verbatim}
claude-o-go -p "Explain this repository in one paragraph."
\end{verbatim}

Switch models for a single invocation:

\begin{verbatim}
OPENCODE_GO_MODEL=minimax-m3 claude-o-go
\end{verbatim}

\section{Why a Proxy Is Needed}

OpenCode Go exposes several models through an Anthropic-compatible Messages API.
It is a routing and billing layer, not a model host. Each request is routed to
the model vendor's own API.

Examples:

\begin{verbatim}
qwen3.7-plus
minimax-m3
qwen3.7-max
\end{verbatim}

Claude Code validates model names locally before sending requests. If you try
to call Claude Code directly with an OpenCode Go model, Claude Code rejects it
before the request is useful:

\begin{verbatim}
ANTHROPIC_BASE_URL=https://opencode.ai/zen/go/v1 \
ANTHROPIC_API_KEY=... \
claude --model qwen3.7-plus
\end{verbatim}

The proxy solves this by letting Claude Code think it is using a normal Claude
alias such as \code{sonnet}, while the proxy rewrites the outgoing request to
use the OpenCode Go model.

\section{Model Compatibility}

Each upstream provider implements its own Anthropic-compatible endpoint, so
compatibility varies. Claude Code sends a full Anthropic-style payload,
including tools, cache control, tool choice, and Anthropic beta headers.

\begin{longtable}{|p{0.24\linewidth}|p{0.22\linewidth}|p{0.43\linewidth}|}
\hline
\textbf{Model id} & \textbf{Provider} & \textbf{Works with Claude Code harness?} \\
\hline
\code{qwen3.7-plus} & Alibaba (Qwen) & Yes \\
\hline
\code{minimax-m3} & MiniMax & Yes \\
\hline
\code{glm-5.2} & Z.ai (Zhipu AI) & No; invalid API parameter \\
\hline
\code{glm-5.1} & Z.ai (Zhipu AI) & No; same as \code{glm-5.2} \\
\hline
\code{kimi-k2.7-code} & Moonshot AI & No; rejects tool function name format \\
\hline
\code{deepseek-v4-pro} & DeepSeek & No; expects OpenAI tool shape \\
\hline
\end{longtable}

Models from the same vendor generally behave like the tested one. Unverified
models should be tested before relying on them.

Practical recommendation:

\begin{verbatim}
export OPENCODE_GO_MODEL="qwen3.7-plus"   # or minimax-m3
\end{verbatim}

\section{Request Flow}

When you run:

\begin{verbatim}
./claude-o-go -p "Reply with exactly: OK"
\end{verbatim}

the flow is:

\begin{enumerate}
  \item \code{claude-o-go} reads the OpenCode Go API key from \code{.env}.
  \item It starts the local proxy at \code{http://127.0.0.1:4141}.
  \item It sets \code{ANTHROPIC\_BASE\_URL=http://127.0.0.1:4141}.
  \item It starts Claude Code with \code{claude --bare --model sonnet}.
  \item Claude Code sends a request to \code{/v1/messages}.
  \item The proxy changes the request body model to the OpenCode Go model.
  \item The proxy forwards the request to \code{https://opencode.ai/zen/go/v1/messages}.
  \item OpenCode Go returns the model response.
  \item The proxy streams the response back to Claude Code.
\end{enumerate}

\section{Important Files}

\subsection{\code{claude-o-go}}

The launcher exports the API key, model, and local base URL, then starts
Claude Code:

\begin{verbatim}
export OPENCODE_GO_API_KEY="$api_key"
export OPENCODE_GO_MODEL="$model"
export ANTHROPIC_BASE_URL="http://127.0.0.1:${proxy_port}"
claude --bare --model "$claude_model_alias" "$@"
\end{verbatim}

Important defaults and environment variables:

\begin{itemize}
  \item \code{OPENCODE\_GO\_MODEL}: required; no hardcoded launcher default.
  \item \code{OPENCODE\_GO\_PROXY\_PORT=4141}.
  \item \code{CLAUDE\_CODE\_MODEL\_ALIAS=sonnet}.
  \item \code{CLAUDE\_CODE\_OPENCODE\_GO\_BARE=1}.
\end{itemize}

\subsection{\code{opencode-go-claude-proxy.mjs}}

The proxy's core rewrite is:

\begin{verbatim}
body.model = model
\end{verbatim}

For generation requests, the proxy streams OpenCode Go's response back to
Claude Code. It handles Node stream backpressure and aborts the upstream
request if Claude Code closes the local connection, so a transient
\code{connection interrupted} message should not take down the proxy process.

The proxy also implements endpoints Claude Code probes:

\begin{verbatim}
HEAD /
HEAD /v1
GET /v1/models
POST /v1/messages/count_tokens
POST /v1/messages
\end{verbatim}

\section{Local Checks}

Before pushing changes, run:

\begin{verbatim}
npm run check
\end{verbatim}

That syntax-checks both JavaScript entry points.

\section{Verification}

Run:

\begin{verbatim}
./claude-o-go -p "Reply with exactly: OK"
\end{verbatim}

Look for proxy log lines like:

\begin{verbatim}
OpenCode Go Claude proxy listening on http://127.0.0.1:4141/v1 -> qwen3.7-plus
POST /v1/messages
POST /messages -> 200 as qwen3.7-plus
OK
\end{verbatim}

The key line is:

\begin{verbatim}
POST /messages -> 200 as qwen3.7-plus
\end{verbatim}

That means Claude Code sent a request to the local proxy, the proxy forwarded it
to OpenCode Go, OpenCode Go returned HTTP 200, and the upstream model was
\code{qwen3.7-plus}.

\section{Common Problems}

\subsection{\code{There's an issue with the selected model}}

Use the launcher:

\begin{verbatim}
./claude-o-go
\end{verbatim}

Do not call \code{claude --model qwen3.7-plus} directly.

\subsection{Double \code{/v1/v1/messages}}

\code{ANTHROPIC\_BASE\_URL} should point to the proxy root:

\begin{verbatim}
http://127.0.0.1:4141
\end{verbatim}

It should not include \code{/v1}, because Claude Code appends
\code{/v1/messages} itself.

\subsection{Claude Code Shows \code{sonnet}}

That is expected. Claude Code sees \code{sonnet}; OpenCode Go receives
\code{qwen3.7-plus}. The proxy log is the source of truth.

\section{Mental Model}

\begin{verbatim}
Claude Code CLI
  thinks model = sonnet
  sends Anthropic request
        |
        v
Local proxy
  rewrites model = qwen3.7-plus
  forwards request
        |
        v
OpenCode Go
  runs the actual model
  returns response
        |
        v
Claude Code CLI
  displays the answer
\end{verbatim}

The Claude Code user experience stays the same, but the model backend is
OpenCode Go.

\end{document}
